\documentclass[11pt,a4paper]{beamer}
\usetheme{Madrid}
\usepackage[utf8]{inputenc}
\usepackage{amsmath}
\usepackage{amsfonts}
\usepackage{amssymb}
\usepackage{graphicx}
\usepackage{tikz}
\usepackage{booktabs}
\usepackage{array}
\usepackage{colortbl}
\usepackage{multicol}

\usefonttheme{professionalfonts}

\setbeamerfont{normal text}{size=\tiny}
\setbeamerfont{itemize/enumerate body}{size=\tiny}
\setbeamerfont{itemize/enumerate subbody}{size=\tiny}
\setbeamerfont{block body}{size=\tiny}
\setbeamerfont{block title}{size=\scriptsize}
\setbeamerfont{frametitle}{size=\small}
\setbeamerfont{title}{size=\small}
\setbeamerfont{subtitle}{size=\footnotesize}
\setbeamerfont{author}{size=\tiny}
\setbeamerfont{date}{size=\tiny}

\setlength{\leftmargini}{0.3em}
\setlength{\leftmarginii}{0.3em}
\setlength{\labelsep}{0.1em}
\setlength{\itemsep}{0.02em}
\setlength{\parskip}{0.02em}

\title[CAMS Exam Preparation]{CAMS Exam Preparation}
\subtitle{AFC Cooperation \& Information Sharing}
\author{FATF Recommendations 36-40 | PPP | FIU Collaboration | Cross-Border}
\date{}

\setbeamercolor{title}{fg=white,bg=blue!70!black}
\setbeamercolor{frametitle}{fg=white,bg=blue!70!black}
\setbeamertemplate{navigation symbols}{}

\begin{document}

% TITLE SLIDE
\begin{frame}
\titlepage
\centering
\tiny 12 Topics | AFC Cooperation | How to Answer
\end{frame}

% ============================================================
% SLIDE 1 - AFC COOPERATION OVERVIEW
% ============================================================
\begin{frame}{1. AFC Cooperation Overview}
\tiny
\noindent\textbf{Introduction to the topic:} Anti-Financial Crime (AFC) cooperation refers to the collaborative efforts between financial institutions, regulators, law enforcement, and international bodies to combat money laundering, terrorist financing, and other illicit finance activities. The core premise is that financial crime is borderless — criminals exploit jurisdictional gaps, and effective countermeasures require information sharing and coordinated action across institutions and national boundaries. Cooperation takes several forms: public-private partnerships where government agencies and financial institutions share intelligence; cross-border collaboration between Financial Intelligence Units (FIUs); and voluntary information sharing between financial institutions, including appropriate foreign financial institutions. The FATF has established specific recommendations (36-40) on international cooperation, requiring countries to provide the widest possible range of mutual legal assistance.

\vspace{0.1cm}
\noindent\textbf{Type of violation or warning signs:} A financial institution that does not cooperate with law enforcement or regulator requests for information. A jurisdiction that refuses to provide mutual legal assistance to other countries. An institution that does not participate in information sharing initiatives with other banks or FIUs. A regulator that does not share information with foreign counterparts. An institution that fails to respond to requests from Financial Intelligence Units in a timely manner. A jurisdiction that has not ratified or implemented international cooperation treaties.

\vspace{0.1cm}
\noindent\textbf{Correct answer approach:} The correct answer will state that AFC cooperation is essential for effective AML/CFT programs, that institutions and jurisdictions must cooperate with law enforcement and regulators, and that information sharing is a key pillar of international AML efforts. Eliminate answers that suggest cooperation is optional or that information sharing violates privacy laws when done for AML purposes. Remember: cooperation includes sharing information with FIU, law enforcement, and other financial institutions, subject to appropriate confidentiality and data protection.
\end{frame}

% ============================================================
% SLIDE 2 - FATF RECOMMENDATIONS 36-40 (INTERNATIONAL COOPERATION)
% ============================================================
\begin{frame}{2. FATF Recommendations 36-40}
\tiny
\noindent\textbf{Introduction to the topic:} FATF Recommendations 36 through 40 address international cooperation in AML/CFT. Recommendation 36 requires countries to provide the widest possible range of mutual legal assistance. Recommendation 37 addresses extradition and mutual legal assistance for money laundering and terrorist financing. Recommendation 38 covers freezing and confiscation of assets across borders. Recommendation 39 addresses extradition. Recommendation 40 requires countries to provide the widest possible range of international cooperation between FIUs and other competent authorities. These recommendations ensure that criminals cannot evade AML enforcement by moving assets or operations across borders.

\vspace{0.1cm}
\noindent\textbf{Type of violation or warning signs:} A country has not ratified mutual legal assistance treaties with key partner countries. A jurisdiction has implemented laws that block international cooperation (e.g., bank secrecy laws that prevent sharing of information with foreign regulators). A country has not designated a Financial Intelligence Unit that can cooperate with foreign FIUs. A jurisdiction has not implemented asset freezing and confiscation provisions that allow cross-border enforcement. A country has extradition restrictions that prevent prosecuting money laundering cases across borders.

\vspace{0.1cm}
\noindent\textbf{Correct answer approach:} The correct answer will state that FATF requires countries to provide mutual legal assistance, cooperate with foreign FIUs, and allow cross-border asset freezing and confiscation. Eliminate answers that suggest a country can opt out of international cooperation or that bank secrecy laws override international obligations. Remember: Recommendations 36-40 are mandatory — countries that do not implement them risk being placed on FATF's Grey List or Black List.
\end{frame}

% ============================================================
% SLIDE 3 - FINANCIAL INTELLIGENCE UNITS (FIU) COOPERATION
% ============================================================
\begin{frame}{3. FIU Cooperation (Egmont Group)}
\tiny
\noindent\textbf{Introduction to the topic:} Financial Intelligence Units are national centers for receiving, analyzing, and disseminating financial intelligence to combat money laundering and terrorist financing. The Egmont Group is the international body that facilitates cooperation between FIUs. It provides a framework for secure information sharing and the Egmont Secure Web, a secure communication channel for FIUs. Under Recommendation 40, FIUs must be able to cooperate and exchange information with foreign FIUs spontaneously and upon request, regardless of the type of partner FIU. Information exchange must be protected by appropriate confidentiality provisions.

\vspace{0.1cm}
\noindent\textbf{Type of violation or warning signs:} A country has not established a FIU that can join the Egmont Group. An FIU refuses to share information with foreign FIUs citing legal obstacles. An FIU has not implemented the Egmont Group principles and standards. An FIU does not have the authority to exchange information spontaneously. An FIU's information sharing is blocked by domestic bank secrecy laws. A financial institution receives requests from foreign FIUs but does not respond or delays response.

\vspace{0.1cm}
\noindent\textbf{Correct answer approach:} The correct answer will state that FIUs must cooperate with foreign FIUs and exchange information spontaneously and upon request. Eliminate answers that suggest FIU cooperation is optional or that FIUs can only share information within their own jurisdiction. Remember: FIU cooperation is essential for global AML/CFT efforts. The Egmont Group provides the framework for this cooperation.
\end{frame}

% ============================================================
% SLIDE 4 - PUBLIC-PRIVATE PARTNERSHIPS (PPP)
% ============================================================
\begin{frame}{4. Public-Private Partnerships (PPP) in AML}
\tiny
\noindent\textbf{Introduction to the topic:} Public-Private Partnerships (PPPs) bring together government agencies (FIUs, law enforcement, regulators) and financial institutions to share information, intelligence, and expertise. PPPs help identify emerging threats, improve suspicious activity reporting, and enable more effective investigations. Examples include the UK's Joint Money Laundering Intelligence Taskforce (JMLIT) and the US's FinCEN Exchange. PPPs are increasingly recognized as essential for effective AML/CFT programs because they leverage private sector intelligence and public sector authority.

\vspace{0.1cm}
\noindent\textbf{Type of violation or warning signs:} A jurisdiction has no PPP framework or mechanism for information sharing between government and financial institutions. Financial institutions do not participate in PPP initiatives in their jurisdiction. Law enforcement does not share feedback on SARs filed by financial institutions. Information sharing is one-way only (institutions report but do not receive intelligence). PPP meetings are infrequent or have low participation. Feedback on the usefulness of reported information is not provided.

\vspace{0.1cm}
\noindent\textbf{Correct answer approach:} The correct answer will state that PPPs are essential for effective AML/CFT and involve sharing information between government and financial institutions, with feedback on SARs being a key element. Eliminate answers that suggest PPPs are optional or that financial institutions should not share information with government. Remember: PPPs improve the quality of SARs and help identify emerging threats. Feedback is critical to maintaining engagement.
\end{frame}

% ============================================================
% SLIDE 5 - CROSS-BORDER INFORMATION SHARING
% ============================================================
\begin{frame}{5. Cross-Border Information Sharing}
\tiny
\noindent\textbf{Introduction to the topic:} Cross-border information sharing involves the exchange of financial intelligence and other information between competent authorities in different countries. This includes FIU-to-FIU sharing under the Egmont Group, law enforcement-to-law enforcement sharing under mutual legal assistance treaties, and regulator-to-regulator sharing. The Wolfsberg Group and other banking associations have issued guidance on information sharing between financial institutions across borders, subject to data protection and confidentiality obligations. Information sharing must be protected by appropriate safeguards to prevent misuse or unauthorized disclosure.

\vspace{0.1cm}
\noindent\textbf{Type of violation or warning signs:} A financial institution refuses to share information with a foreign regulator on the basis of bank secrecy laws. A jurisdiction has laws that prohibit the sharing of financial information with foreign authorities. An FIU does not have a mechanism for exchanging information with foreign FIUs. A country has not signed mutual legal assistance treaties with key partner countries. A financial institution does not have procedures for responding to cross-border information requests. Information sharing agreements are not documented or are not respected.

\vspace{0.1cm}
\noindent\textbf{Correct answer approach:} The correct answer will state that cross-border information sharing is essential and must be facilitated by FIUs, law enforcement, and regulators, with appropriate confidentiality protections. Eliminate answers that suggest bank secrecy laws override the obligation to share information for AML purposes. Remember: FATF Recommendation 40 requires countries to provide the widest possible range of international cooperation.
\end{frame}

% ============================================================
% SLIDE 6 - INFORMATION SHARING BETWEEN FINANCIAL INSTITUTIONS
% ============================================================
\begin{frame}{6. Information Sharing Between Financial Institutions}
\tiny
\noindent\textbf{Introduction to the topic:} Information sharing between financial institutions allows banks to share intelligence about suspicious customers, transactions, or patterns, particularly in correspondent banking relationships or multi-bank money laundering schemes. Many jurisdictions have enacted laws that provide safe harbor for sharing information for AML purposes, including sharing SAR-related information. However, financial institutions must be careful to avoid tipping off customers who are the subject of a SAR. Information sharing should be governed by written agreements that specify the purpose and confidentiality of shared information.

\vspace{0.1cm}
\noindent\textbf{Type of violation or warning signs:} A financial institution refuses to share information with another financial institution that is also affected by the same suspicious activity, citing privacy or confidentiality concerns. A financial institution shares information without appropriate confidentiality safeguards. A financial institution shares information that would tip off the customer that they are under investigation. Information sharing agreements are not in place or are not followed. A financial institution shares information for non-AML purposes (e.g., for commercial gain) without authorization.

\vspace{0.1cm}
\noindent\textbf{Correct answer approach:} The correct answer will state that financial institutions can and should share information for AML purposes, subject to appropriate confidentiality and safe harbor provisions, and must not tip off customers. Eliminate answers that suggest financial institutions should never share information or that sharing information always violates privacy laws. Remember: safe harbor protects institutions that share information in good faith for AML purposes.
\end{frame}

% ============================================================
% SLIDE 7 - MUTUAL LEGAL ASSISTANCE (MLA)
% ============================================================
\begin{frame}{7. Mutual Legal Assistance (MLA)}
\tiny
\noindent\textbf{Introduction to the topic:} Mutual Legal Assistance (MLA) is the formal process by which countries request and provide assistance in criminal investigations, including money laundering and terrorist financing cases. MLA requests can include seizure of assets, freezing of accounts, taking of evidence, and extradition. Under FATF Recommendation 36, countries must provide the widest possible range of mutual legal assistance in relation to money laundering and terrorist financing offenses. MLA must be timely and effective, and countries must not refuse assistance solely on grounds of bank secrecy.

\vspace{0.1cm}
\noindent\textbf{Type of violation or warning signs:} A country refuses to provide MLA to another country on the grounds of bank secrecy or tax offenses. A country has no MLA treaty with key partner countries. A country's MLA process is excessively slow or bureaucratic. A country refuses to freeze or confiscate assets that are the subject of an MLA request. A country does not have procedures for executing MLA requests related to money laundering or terrorist financing. A country does not provide MLA for cases involving foreign PEPs.

\vspace{0.1cm}
\noindent\textbf{Correct answer approach:} The correct answer will state that countries must provide MLA for money laundering and terrorist financing offenses and cannot refuse on the basis of bank secrecy. Eliminate answers that suggest MLA is optional or that tax offenses are not covered. Remember: MLA is a cornerstone of international cooperation — countries that fail to provide MLA risk being listed as non-cooperative.
\end{frame}

% ============================================================
% SLIDE 8 - ASSET FREEZING AND CONFISCATION ACROSS BORDERS
% ============================================================
\begin{frame}{8. Asset Freezing and Confiscation Across Borders}
\tiny
\noindent\textbf{Introduction to the topic:} Freezing and confiscation of assets across borders is essential to prevent criminals from moving illicit funds to jurisdictions with weak enforcement. Under FATF Recommendation 38, countries must have the authority to take swift action in response to requests from another country for the freezing or confiscation of property associated with money laundering or terrorist financing. This includes the ability to freeze assets without prior notification to the owner. Countries must also have procedures for recognizing and enforcing foreign confiscation orders.

\vspace{0.1cm}
\noindent\textbf{Type of violation or warning signs:} A country refuses to freeze assets requested by a foreign jurisdiction. A country has no mechanism for recognizing foreign confiscation orders. A country's asset freezing process is slow or bureaucratic. A country tips off the asset owner before executing a freezing request. A country does not have procedures for confiscating assets that are identified as proceeds of crime by a foreign jurisdiction. A country allows assets to be moved out before freezing takes effect.

\vspace{0.1cm}
\noindent\textbf{Correct answer approach:} The correct answer will state that countries must have authority to freeze and confiscate assets at the request of foreign jurisdictions, without prior notification to the owner. Eliminate answers that suggest asset freezing is only available domestically or that prior notification is required. Remember: speed is essential — criminals will move assets if they are tipped off.
\end{frame}

% ============================================================
% SLIDE 9 - COOPERATION IN CORRESPONDENT BANKING
% ============================================================
\begin{frame}{9. Cooperation in Correspondent Banking}
\tiny
\noindent\textbf{Introduction to the topic:} Correspondent banking relationships require significant cooperation between the correspondent bank (which provides the account) and the respondent bank (which uses the account). FATF Recommendation 13 requires banks to obtain sufficient information about the respondent bank, understand its business and AML controls, and obtain senior management approval. The respondent bank must cooperate by providing information about its own AML program, its clients (including PEPs), and its transaction monitoring. Failure of cooperation can lead to termination of the correspondent relationship.

\vspace{0.1cm}
\noindent\textbf{Type of violation or warning signs:} A respondent bank refuses to provide information about its AML program to the correspondent bank. A respondent bank refuses to identify its underlying clients (including PEPs) to the correspondent bank. A respondent bank does not provide timely responses to due diligence requests. A correspondent bank continues the relationship despite repeated failures of cooperation. A respondent bank has no independent AML audit and refuses to obtain one. A respondent bank operates in a jurisdiction with bank secrecy laws that prevent cooperation.

\vspace{0.1cm}
\noindent\textbf{Correct answer approach:} The correct answer will state that correspondent banking requires full cooperation and information sharing, and failure to cooperate is grounds for termination. Eliminate answers that suggest the correspondent bank has no obligation to seek information or that the respondent bank's bank secrecy laws excuse non-cooperation. Remember: without cooperation, the correspondent bank is blind to the respondent bank's clients and risks.
\end{frame}

% ============================================================
% SLIDE 10 - EXTRADITION AND CROSS-BORDER PROSECUTION
% ============================================================
\begin{frame}{10. Extradition and Cross-Border Prosecution}
\tiny
\noindent\textbf{Introduction to the topic:} Extradition allows one country to request the transfer of a criminal suspect from another country for prosecution. Under FATF Recommendation 39, countries must have extradition arrangements that allow for the extradition of individuals charged with money laundering and terrorist financing offenses. Countries must not refuse extradition solely on the grounds that the offense involves fiscal matters (i.e., tax offenses) or that the offense is not an offense in the requested country. Extradition should be timely and efficient.

\vspace{0.1cm}
\noindent\textbf{Type of violation or warning signs:} A country refuses to extradite a money laundering suspect to another country on the grounds of tax offenses. A country has no extradition treaty with a key partner country. A country's extradition process is excessively slow. A country refuses to extradite its own nationals for money laundering offenses. A country provides a safe haven for individuals wanted for money laundering by refusing to extradite. A country does not have procedures for extradition in money laundering cases.

\vspace{0.1cm}
\noindent\textbf{Correct answer approach:} The correct answer will state that countries must extradite individuals charged with money laundering offenses and cannot refuse on the grounds of tax offenses or fiscal matters. Eliminate answers that suggest extradition is optional or that nationals cannot be extradited. Remember: countries that refuse extradition risk being seen as safe havens for money launderers.
\end{frame}

% ============================================================
% SLIDE 11 - FIU SPONTANEOUS INFORMATION EXCHANGE
% ============================================================
\begin{frame}{11. FIU Spontaneous Information Exchange}
\tiny
\noindent\textbf{Introduction to the topic:} Spontaneous information exchange occurs when an FIU proactively shares information with a foreign FIU without waiting for a request. This is particularly important in counter-terrorist financing and other urgent cases where time is critical. Under the Egmont Group principles, FIUs are encouraged to share information spontaneously when it is relevant to a foreign FIU's work. Spontaneous exchange helps identify cross-border connections that might otherwise be missed. The information must be protected by appropriate confidentiality provisions.

\vspace{0.1cm}
\noindent\textbf{Type of violation or warning signs:} An FIU does not have authority to share information spontaneously. An FIU never shares information spontaneously, only in response to requests. An FIU shares information spontaneously but without adequate confidentiality safeguards. An FIU does not provide feedback on information received from foreign FIUs. An FIU's spontaneous sharing is blocked by domestic legal obstacles. An FIU does not have procedures for spontaneous information exchange.

\vspace{0.1cm}
\noindent\textbf{Correct answer approach:} The correct answer will state that FIUs should share information spontaneously when relevant to a foreign FIU's work, subject to confidentiality safeguards. Eliminate answers that suggest spontaneous sharing is not permitted or only occurs on request. Remember: spontaneous sharing can help identify terrorist financing and other urgent cross-border threats before they materialize.
\end{frame}

% ============================================================
% SLIDE 12 - CHALLENGES TO AFC COOPERATION
% ============================================================
\begin{frame}{12. Challenges to AFC Cooperation}
\tiny
\noindent\textbf{Introduction to the topic:} AFC cooperation faces several challenges: differing legal standards and definitions of money laundering across jurisdictions; bank secrecy laws that prevent sharing of customer information; concerns about data protection and privacy; lack of resources and capacity in smaller jurisdictions; political obstacles to cooperation; and concerns about the use of shared information by law enforcement. Overcoming these challenges requires trust building, capacity building, and clear legal frameworks that allow for information sharing while protecting privacy rights.

\vspace{0.1cm}
\noindent\textbf{Type of violation or warning signs:} A jurisdiction has bank secrecy laws that prevent sharing of information. A jurisdiction has weak data protection laws that undermine trust in sharing. A jurisdiction has not implemented FATF recommendations on international cooperation. A jurisdiction's FIU lacks resources to respond to foreign requests. A jurisdiction refuses to cooperate with another country for political reasons. Information sharing agreements are not respected by one party. Different definitions of money laundering across jurisdictions create legal barriers.

\vspace{0.1cm}
\noindent\textbf{Correct answer approach:} The correct answer will state that challenges include legal obstacles, data protection concerns, and lack of resources, but these must be overcome through legal reform, capacity building, and trust. Eliminate answers that suggest challenges are insurmountable or that cooperation should not be attempted. Remember: cooperation is essential despite challenges, and FATF expects countries to address these barriers.
\end{frame}

% ============================================================
% SUMMARY SLIDE - AFC COOPERATION PRINCIPLES
% ============================================================
\begin{frame}{Summary: Key Principles for AFC Cooperation}
\tiny
\noindent\textbf{How to answer exam questions on AFC cooperation:}

\vspace{0.1cm}
\noindent• \textbf{Mutual Legal Assistance (MLA):} Countries must provide MLA for ML/TF offenses and cannot refuse on grounds of bank secrecy or tax offenses. Required under FATF Rec. 36.

\vspace{0.1cm}
\noindent• \textbf{FIU Cooperation (Egmont Group):} FIUs must exchange information with foreign FIUs both upon request and spontaneously. Required under FATF Rec. 40.

\vspace{0.1cm}
\noindent• \textbf{Asset Freezing/Confiscation:} Countries must freeze and confiscate assets at the request of foreign jurisdictions without prior notification. Required under FATF Rec. 38.

\vspace{0.1cm}
\noindent• \textbf{Extradition:} Countries must extradite individuals charged with ML/TF offenses and cannot refuse on fiscal grounds. Required under FATF Rec. 39.

\vspace{0.1cm}
\noindent• \textbf{Public-Private Partnerships:} PPPs enable information sharing between government and financial institutions, improving SAR quality.

\vspace{0.1cm}
\noindent• \textbf{Information Sharing Between Institutions:} Financial institutions can share information for AML purposes, with safe harbor and without tipping off.

\vspace{0.1cm}
\noindent• \textbf{Correspondent Banking:} Respondent banks must cooperate and share information with correspondent banks. Failure = termination.

\vspace{0.1cm}
\noindent• \textbf{Challenges:} Legal obstacles, data protection, and resource constraints can be overcome through trust-building and legal reform.

\vspace{0.1cm}
\noindent• \textbf{Grey/Black List Risk:} Countries that do not cooperate internationally risk being placed on FATF's Grey or Black List.

\vspace{0.3cm}
\centering
{Remember: Cooperation is mandatory, not optional. When in doubt, share information (subject to confidentiality) and cooperate with regulators and FIUs.}
\end{frame}

\end{document}